\begin{figure*}[!t]
\centering
\fbox{\begin{minipage}{0.95\textwidth}
\textbf{Algorithm 3: Rejection-Aware Residual Resource Allocation}\\[0.35em]
\textbf{Input:} edge-selected tasks $\mathcal{I}_e$, task/resource features, demand-proportional base allocation $\mathbf{b}_i$, trained residual network $f_{\phi}$.\\
\textbf{Output:} capacity-feasible allocation $\hat{\mathbf{y}}_i$ for CPU, memory, and bandwidth.\\[0.35em]
1) For each edge-selected task, compute the demand-proportional base allocation $\mathbf{b}_i$.\\
2) Concatenate task descriptors, edge state, network state, policy outputs, and the bounded pre-decision QoS/admission-risk estimate $\hat{\rho}_i$ into allocator input $z_i$; do not include realized labels $M_i$, $R_i$, $Q_i$, Oracle actions, or generated final allocation targets.\\
3) Predict the residual correction $f_{\phi}(z_i)$ and form $\tilde{\mathbf{y}}_i=\mathbf{b}_i+f_{\phi}(z_i)$.\\
4) Build the rejection-aware demand branch from urgency, deadline, priority, edge-SLA risk, and related pre-decision risk features; blend it with the capacity-normalized residual branch.\\
5) Project the result into valid normalized capacity so CPU, memory, and bandwidth assignments remain feasible.\\
6) Report MAE, under-allocation, waste, capacity violation, estimated SLA risk, and efficiency.
\end{minipage}}
\caption{Stage-2 allocation procedure for tasks already routed to edge execution by \fedddqn.}
\label{alg:stage2_allocator}
\end{figure*}
